% Options for packages loaded elsewhere
\PassOptionsToPackage{unicode}{hyperref}
\PassOptionsToPackage{hyphens}{url}
\documentclass[
  a4paper,11pt]{article}
\usepackage{xcolor}
\usepackage[margin=1in]{geometry}
\usepackage{amsmath,amssymb}
\usepackage{graphicx}
\setcounter{secnumdepth}{-\maxdimen} % remove section numbering
\usepackage{iftex}
\ifPDFTeX
  \usepackage[T1]{fontenc}
  \usepackage[utf8]{inputenc}
  \usepackage{textcomp}
\else
  \usepackage{unicode-math}
  \defaultfontfeatures{Scale=MatchLowercase}
  \defaultfontfeatures[\rmfamily]{Ligatures=TeX,Scale=1}
\fi
\usepackage{lmodern}
\ifPDFTeX\else
  \setmainfont[]{Times New Roman}
\fi
\IfFileExists{upquote.sty}{\usepackage{upquote}}{}
\IfFileExists{microtype.sty}{%
  \usepackage[]{microtype}
  \UseMicrotypeSet[protrusion]{basicmath}
}{}
\makeatletter
\@ifundefined{KOMAClassName}{%
  \IfFileExists{parskip.sty}{%
    \usepackage{parskip}
  }{%
    \setlength{\parindent}{0pt}
    \setlength{\parskip}{6pt plus 2pt minus 1pt}}
}{% if KOMA class
  \KOMAoptions{parskip=half}}
\makeatother
\usepackage{longtable,booktabs,array}
\newcounter{none}
\usepackage{calc}
\usepackage{etoolbox}
\makeatletter
\patchcmd\longtable{\par}{\if@noskipsec\mbox{}\fi\par}{}{}
\makeatother
\IfFileExists{footnotehyper.sty}{\usepackage{footnotehyper}}{\usepackage{footnote}}
\makesavenoteenv{longtable}
\usepackage{bookmark}
\IfFileExists{xurl.sty}{\usepackage{xurl}}{}
\urlstyle{same}
\usepackage{hyperref}
\hypersetup{
  hidelinks,
  pdfcreator={LaTeX via pandoc}}
\setlength{\emergencystretch}{3em}
\providecommand{\tightlist}{%
  \setlength{\itemsep}{0pt}\setlength{\parskip}{0pt}}
\usepackage{amsthm}
\newtheorem{theorem}{Theorem}[section]
\newtheorem{proposition}[theorem]{Proposition}
\newtheorem{lemma}[theorem]{Lemma}
\newtheorem{corollary}[theorem]{Corollary}
\theoremstyle{definition}
\newtheorem{definition}[theorem]{Definition}
\theoremstyle{remark}
\newtheorem{remark}[theorem]{Remark}
\newtheorem{observation}[theorem]{Observation}
\begin{document}
\section{Four Modes of Shape-Invariant Waves: Spin, Chirality, and
Statistics Determined by Wave Vector
Structure}\label{four-modes-of-shape-invariant-waves-spin-chirality-and-statistics-determined-by-wave-vector-structure}

\textbf{Author}: Noriaki Kihara\\
\textbf{Affiliation}: WF System Co., Ltd.\\
\textbf{ORCID}: 0009-0004-6753-4020\\
\textbf{Version}: v1.0\\
\textbf{Date}: April 2026\\
\textbf{DOI}: 10.5281/zenodo.19709798

\begin{center}\rule{0.5\linewidth}{0.5pt}\end{center}

\subsection{Abstract}\label{abstract}

Building on the shape-invariant standing waves on the closed sphere
\(S^3\) introduced in the companion paper {[}1{]} and the
shape-invariant traveling waves on \(\mathbb{Z}^4\) formulated in
{[}2{]}, we classify shape-invariant waves based on the number of
nonzero components \(\kappa\) of the wave vector. Waves with
\(\kappa = 2\) (two nonzero components) constitute direction-constrained
waves whose displacement direction is restricted to one spatial
dimension, and the signed area of their two-axis plane admits only two
states, thereby providing a geometric origin for the Pauli exclusion
principle. Waves with \(\kappa = 0\) (all spatial components zero)
propagate spherically; when the temporal component \(k_t = 0\), they are
time-reversal symmetric and thus carry no intrinsic causal ordering.
Waves with \(\kappa = 1\) (only one spatial component nonzero, temporal
component zero) are standing waves that possess scalar character due to
the equivalence of three directions and fill the entire subjective space
as a ground mode. Direction-constrained waves with \(\kappa = 2\)
possessing a (space, \(t\)) plane have chirality defined by the
orientation of the \(t\)-axis, whereas waves possessing only (space,
space) planes have degenerate chirality and may satisfy the Majorana
condition. Central projection onto the child space induces geodesic
curvature, and acceleration is expressed as a discrete update rule for
the wave vector.

\textbf{Keywords}: shape-invariant wave, mode classification,
direction-constrained wave, signed area, chirality, Pauli exclusion
principle, Majorana condition, central projection, acceleration

\begin{center}\rule{0.5\linewidth}{0.5pt}\end{center}

\subsection{§1 Introduction and Scope of
Claims}\label{introduction-and-scope-of-claims}

This paper demonstrates, within the framework of shape-invariant waves
constructed in the companion papers {[}1{]} and {[}2{]}, that the
structure of the wave vector determines four qualitatively distinct
modes, and that the geometric properties of each mode provide the
origins of spin, chirality, and statistics.

The claims of this paper are limited to the following five points:

\begin{enumerate}
\def\labelenumi{\arabic{enumi}.}
\tightlist
\item
  The number of nonzero components \(\kappa\) of the wave vector
  classifies shape-invariant waves into three types:
  \(\kappa \in \{0, 1, 2\}\). No stable modes exist for
  \(\kappa \geq 3\).
\item
  Waves with \(\kappa = 2\) admit only two states of signed area
  orientation in their two-axis plane, geometrically prohibiting the
  coexistence of three or more waves with the same wave vector.
\item
  For waves with \(\kappa = 2\), the (space, \(t\)) plane defines
  chirality from the fixed direction of the \(t\)-axis, whereas
  chirality is degenerate in the (space, space) plane.
\item
  Spherical waves with \(\kappa = 0\) and \(k_t = 0\) carry no causal
  ordering, and the shortest geodesic path is retroactively constructed
  from the interaction outcome.
\item
  Central projection onto the child space induces geodesic curvature on
  wave trajectories and geometrically expresses acceleration.
\end{enumerate}

The following elements are \textbf{outside the scope of this paper}:

\begin{itemize}
\tightlist
\item
  Identification with specific physical particles
\item
  Numerical derivation of coupling constants
\item
  Reproduction in the continuous limit of gauge groups
\item
  Numerical derivation of mass values and mixing matrices
\end{itemize}

\textbf{Note on coordinate systems}: This paper primarily works on
\(\mathbb{Z}^4\) (four axes \(xyzt\)). However, to avoid the problem
that spherical waves with \(\kappa_s = 0\) become trivial (constant
functions) within \(\mathbb{Z}^4\), the extended axes introduced in
{[}4{]} (\(R\): radial direction of central projection, \(Q\): discrete
label axis) are invoked where necessary. See {[}4{]} for details on
these extended axes.

\begin{center}\rule{0.5\linewidth}{0.5pt}\end{center}

\subsection{§2 3+1 Decomposition of the Wave Vector and Mode
Classification}\label{decomposition-of-the-wave-vector-and-mode-classification}

\subsubsection{§2.1 3+1 Decomposition}\label{decomposition}

As shown in §11 of {[}2{]}, the existence condition for shape-invariant
traveling waves and the stability condition for standing waves on a
closed sphere point to the structure of de Sitter spacetime
\(\mathrm{dS}_4\). Within this structure, the four axes of
\(\mathbb{Z}^4\) naturally decompose into 3+1:

\begin{itemize}
\tightlist
\item
  \textbf{Spatial components}:
  \(\mathbf{k}_s = (k_1, k_2, k_3) \in \mathbb{Z}^3\) (corresponding to
  the three directions of \(S^3\))
\item
  \textbf{Temporal component}: \(k_t \in \mathbb{Z}\) (corresponding to
  the hyperbolic time direction)
\end{itemize}

\subsubsection{\texorpdfstring{§2.2 Number of Nonzero Components
\(\kappa\)}{§2.2 Number of Nonzero Components \textbackslash kappa}}\label{number-of-nonzero-components-kappa}

For a wave vector \(\mathbf{k} = (k_1, k_2, k_3, k_t)\), we define:

\[\kappa_s = |\{i \in \{1, 2, 3\} : k_i \neq 0\}| \quad (\text{number of nonzero spatial components})\]

\[\kappa_t = \begin{cases} 1 & (k_t \neq 0) \\ 0 & (k_t = 0) \end{cases} \quad (\text{temporal nonzero indicator})\]

\[\kappa = \kappa_s + \kappa_t \quad (\text{total number of nonzero components})\]

The possible values of \(\kappa\) are \(\{0, 1, 2, 3, 4\}\), but by
§2.3, \(\kappa \geq 3\) admits no stable modes.

\subsubsection{\texorpdfstring{§2.3 Non-Existence of
\(\kappa \geq 3\)}{§2.3 Non-Existence of \textbackslash kappa \textbackslash geq 3}}\label{non-existence-of-kappa-geq-3}

Waves with \(\kappa \geq 3\) have simultaneous phase variations in three
or more directions. The following two \textbf{independently sufficient}
reasons show that they do not form stable modes (either one alone
excludes \(\kappa \geq 3\)):

\textbf{(i) Conflict with the Nyquist stability condition} (sufficient
on its own): By Proposition 9.1 of {[}2{]}, the phase change along each
axis is restricted to \(|\phi_i| \leq \pi/2\). The total phase load of a
self-returning loop maintaining nonzero phase changes simultaneously on
\(\kappa\) axes is at least \(\kappa \cdot |\phi_{\min}|\), and for
\(\kappa = 3\), \(3 \cdot |\phi_{\min}| > \pi\) (exceeding the stability
boundary \(\pi\)). Therefore, for \(\kappa \geq 3\), self-returning
loops cannot form stable closed paths.

\textbf{(ii) Directional delocalization} (sufficient on its own): Waves
with \(\kappa_s \geq 2\) and \(\kappa_t = 1\) propagate simultaneously
in multiple spatial directions. They cannot concentrate displacement in
a specific spatial region while satisfying the shape-invariance
condition \(|\psi(\mathbf{x} + \mathbf{a})| = |\psi(\mathbf{x})|\) in
all directions, becoming indistinguishable from the background field.

\textbf{Proposition 2.1}: Shape-invariant traveling waves on
\(\mathbb{Z}^4\) with \(\kappa \geq 3\) have no stable modes. Conditions
(i) and (ii) are each independently sufficient to establish this. The
stable modes of shape-invariant waves are limited to three types:
\(\kappa \in \{0, 1, 2\}\).

\subsubsection{§2.4 Mode Classification
Table}\label{mode-classification-table}

{\def\LTcaptype{none} % do not increment counter
\begin{longtable}[]{@{}
  >{\raggedright\arraybackslash}p{(\linewidth - 8\tabcolsep) * \real{0.0984}}
  >{\raggedright\arraybackslash}p{(\linewidth - 8\tabcolsep) * \real{0.1639}}
  >{\raggedright\arraybackslash}p{(\linewidth - 8\tabcolsep) * \real{0.1967}}
  >{\raggedright\arraybackslash}p{(\linewidth - 8\tabcolsep) * \real{0.1967}}
  >{\raggedright\arraybackslash}p{(\linewidth - 8\tabcolsep) * \real{0.3443}}@{}}
\toprule\noalign{}
\begin{minipage}[b]{\linewidth}\raggedright
Mode
\end{minipage} & \begin{minipage}[b]{\linewidth}\raggedright
\(\kappa\)
\end{minipage} & \begin{minipage}[b]{\linewidth}\raggedright
\(\kappa_s\)
\end{minipage} & \begin{minipage}[b]{\linewidth}\raggedright
\(\kappa_t\)
\end{minipage} & \begin{minipage}[b]{\linewidth}\raggedright
Geometric character
\end{minipage} \\
\midrule\noalign{}
\endhead
\bottomrule\noalign{}
\endlastfoot
Direction-constrained wave & 2 & 1 & 1 & Linear propagation in the (one
spatial axis, \(t\)) plane \\
Planar standing wave \({}^{*}\) & 2 & 2 & 0 & Standing pattern in the
(space, space) plane \\
Spherical wave & 0 & 0 & 0 or 1 & No spatial direction; isotropic \\
Standing wave & 1 & 1 & 0 & Standing oscillation along one spatial axis;
scalar \\
\end{longtable}
}

\({}^{*}\) The planar standing wave shares \(\kappa = 2\) and therefore
shares the signed-area property (exclusion principle) of §7 with the
direction-constrained wave. However, having \(\kappa_t = 0\), it is a
standing structure without propagation and is qualitatively different
from the direction-constrained wave. Details are given in §3.4.

\begin{center}\rule{0.5\linewidth}{0.5pt}\end{center}

\subsection{\texorpdfstring{§3 Direction-Constrained Waves
(\(\kappa = 2\), \(\kappa_s = 1\),
\(\kappa_t = 1\))}{§3 Direction-Constrained Waves (\textbackslash kappa = 2, \textbackslash kappa\_s = 1, \textbackslash kappa\_t = 1)}}\label{direction-constrained-waves-kappa-2-kappa_s-1-kappa_t-1}

\subsubsection{§3.1 Linear Propagation}\label{linear-propagation}

For \(\kappa = 2\) with \(\kappa_s = 1\), \(\kappa_t = 1\), the wave
vector can be written as \(\mathbf{k} = (k_x, 0, 0, k_t)\) (the choice
of spatial axis is without loss of generality by the \(\mathrm{SO}(3)\)
symmetry of \(S^3\)). From the shape-invariance condition:

\[\psi(\mathbf{x} + \mathbf{e}_1) = e^{ik_x} \cdot \psi(\mathbf{x}) \quad (\text{phase change in the $x$ direction})\]

\[\psi(\mathbf{x} + \mathbf{e}_2) = \psi(\mathbf{x}), \quad \psi(\mathbf{x} + \mathbf{e}_3) = \psi(\mathbf{x}) \quad (\text{invariant in the $y, z$ directions})\]

\[\psi(\mathbf{x} + \mathbf{e}_4) = e^{ik_t} \cdot \psi(\mathbf{x}) \quad (\text{phase change in the $t$ direction})\]

This wave has spatial displacement only in the \(x\) direction and no
displacement in the \(y, z\) directions.

\textbf{Proposition 3.1}: A shape-invariant traveling wave with
\(\kappa = 2\), \(\kappa_t = 1\) is a direction-constrained wave that
propagates in only one of the three spatial directions. Its displacement
is constrained to one spatial dimension.

\subsubsection{§3.2 Propagation Velocity and Displacement
Constraint}\label{propagation-velocity-and-displacement-constraint}

The propagation velocity is determined by \(v = k_x / k_t\). By
Proposition 9.1 of {[}2{]} (Nyquist stability condition), \(v \leq c\).

{\def\LTcaptype{none} % do not increment counter
\begin{longtable}[]{@{}lll@{}}
\toprule\noalign{}
Condition & Propagation velocity & Nature of displacement \\
\midrule\noalign{}
\endhead
\bottomrule\noalign{}
\endlastfoot
\(|k_x| = |k_t|\) & \(v = c\) (maximum) & Light-speed propagation \\
\(|k_x| < |k_t|\) & \(v < c\) (subluminal) & Massive \\
\(k_x = 0\) & \(v = 0\) & Degenerates to \(\kappa = 1\) (standing) \\
\end{longtable}
}

The case \(v < c\) occurs when \(|k_x / k_t| < 1\), i.e., the spatial
displacement per temporal component is limited. This displacement
constraint is intrinsic as a component ratio of the wave vector itself,
not an external parameter.

\subsubsection{§3.3 Generation Structure}\label{generation-structure}

Depending on which of the spatial axes \(k_1, k_2, k_3\) is nonzero,
three equivalent but distinguishable direction-constrained waves exist:

{\def\LTcaptype{none} % do not increment counter
\begin{longtable}[]{@{}ll@{}}
\toprule\noalign{}
Nonzero spatial axis & Wave vector \\
\midrule\noalign{}
\endhead
\bottomrule\noalign{}
\endlastfoot
\(k_1\) & \((k_1, 0, 0, k_t)\) \\
\(k_2\) & \((0, k_2, 0, k_t)\) \\
\(k_3\) & \((0, 0, k_3, k_t)\) \\
\end{longtable}
}

By the \(\mathrm{SO}(3)\) symmetry of \(S^3\), these are locally
equivalent but are distinguished as wave vector labels. In Table A of
{[}4{]}, these three choices correspond to the generation structure. The
following naming convention of ``first generation, second generation,
third generation'' indicates correspondence with the interpretation
layer (§9) of {[}4{]} and does not affect the validity of the
propositions in this paper.

\subsubsection{\texorpdfstring{§3.4 Degenerate Type (\(\kappa_s = 2\),
\(\kappa_t = 0\))}{§3.4 Degenerate Type (\textbackslash kappa\_s = 2, \textbackslash kappa\_t = 0)}}\label{degenerate-type-kappa_s-2-kappa_t-0}

For \(\kappa = 2\) with \(\kappa_s = 2\), \(\kappa_t = 0\), the wave
vector is \(\mathbf{k} = (k_x, k_y, 0, 0)\). Since the temporal
component is zero:

\[\psi(\mathbf{x}, t+1) = \psi(\mathbf{x}, t) \quad (\text{time-invariant})\]

This wave is a standing pattern in space without temporal evolution and
is qualitatively different from the direction-constrained wave of §3.1.
However, since \(\kappa = 2\), it shares the signed-area properties
shown in §7.

In Table A of {[}4{]}, the degenerate type corresponds to neutrinos
(\(\nu_e, \nu_\mu, \nu_\tau\)). Each generation's neutrino shares one
spatial axis with the charged fermion of the same generation while
placing the other nonzero axis on a different spatial axis rather than
the time axis:

{\def\LTcaptype{none} % do not increment counter
\begin{longtable}[]{@{}lllll@{}}
\toprule\noalign{}
Generation & Charged fermion & Shared axis & Neutrino & Proper axis \\
\midrule\noalign{}
\endhead
\bottomrule\noalign{}
\endlastfoot
1 & \((k_1, 0, 0, k_t)\) & \(k_1\) & \((k_1, k_2, 0, 0)\) & \(k_2\) \\
2 & \((0, k_2, 0, k_t)\) & \(k_2\) & \((0, k_2, k_3, 0)\) & \(k_3\) \\
3 & \((0, 0, k_3, k_t)\) & \(k_3\) & \((k_1, 0, k_3, 0)\) & \(k_1\) \\
\end{longtable}
}

\begin{center}\rule{0.5\linewidth}{0.5pt}\end{center}

\subsection{\texorpdfstring{§4 Spherical Waves
(\(\kappa_s = 0\))}{§4 Spherical Waves (\textbackslash kappa\_s = 0)}}\label{spherical-waves-kappa_s-0}

\subsubsection{§4.1 Absence of Spatial
Directionality}\label{absence-of-spatial-directionality}

Waves with \(\kappa_s = 0\) have all spatial components equal to zero:
\(\mathbf{k}_s = (0, 0, 0)\). From the shape-invariance condition:

\[\psi(\mathbf{x} + \mathbf{e}_s) = \psi(\mathbf{x}) \quad (\text{phase-invariant in all spatial directions},\; s \in \{1, 2, 3\})\]

A wave with no phase variation in any spatial direction has no preferred
direction in space. When emitted from a point source, it propagates as
an isotropic spherical wave expanding in all directions.

\textbf{Proposition 4.1}: Waves with \(\kappa_s = 0\) are isotropic in
space and propagate as spherical waves in all directions.

\subsubsection{\texorpdfstring{§4.2 Representation on
\(\mathbb{Z}^4\)}{§4.2 Representation on \textbackslash mathbb\{Z\}\^{}4}}\label{representation-on-mathbbz4}

Within the framework of \(\mathbb{Z}^4\), \(\kappa = 0\)
(\(\kappa_s = 0\) and \(\kappa_t = 0\)) reduces to a trivial constant
function. Non-trivial spherical waves are realized by having nonzero
components on the extended axes introduced in {[}4{]} (\(R\): radial
direction of central projection, \(Q\): discrete label axis).

Specifically, a spherical wave with \(\kappa_t = 1\) has
\(\mathbf{k} = (0, 0, 0, k_t)\) and possesses phase variation only in
the temporal direction. It is spatially isotropic but undergoes temporal
evolution.

\subsubsection{§4.3 Time Symmetry and Retroactive Construction of
Causality}\label{time-symmetry-and-retroactive-construction-of-causality}

A spherical wave with \(k_t = 0\) is phase-invariant in the temporal
direction as well:

\[\psi(\mathbf{x}, t + 1) = \psi(\mathbf{x}, t) \quad (\text{time-reversal symmetric})\]

\textbf{Proposition 4.2}: Waves with \(k_t = 0\) are time-reversal
symmetric and carry no intrinsic causal ordering of propagation.

As a consequence of this proposition, the interaction between two
direction-constrained waves mediated by a spherical wave with
\(k_t = 0\) has the following properties:

\begin{enumerate}
\def\labelenumi{\arabic{enumi}.}
\tightlist
\item
  The temporal order of emission and absorption is undefined.
\item
  Only the outcome of the interaction is physically determined.
\item
  From the outcome, the propagation path along the shortest geodesic is
  retroactively constructed.
\end{enumerate}

This provides a geometric derivation of the principle of least action:
the spherical wave does not ``choose'' the shortest path; rather, the
shortest path is retroactively determined from the outcome of an
interaction mediated by a causally unordered spherical wave.

\subsubsection{§4.4 Filling of the Subjective Space as a Ground
Mode}\label{filling-of-the-subjective-space-as-a-ground-mode}

Based on the standing wave condition
\(2\pi R = n \cdot \lambda_{\text{base}}\) from {[}1{]}, when a
spherical wave takes the ground mode (\(n = 1\)), a wave with wavelength
\(\lambda = 2\pi R\) extends over the entire \(S^3\).

The ground-mode spherical wave with \(\kappa_s = 0\): - Uniformly fills
the entire \(S^3\) - Has no phase structure in any spatial direction
(isotropic) - Is temporally invariant (when \(k_t = 0\)) or oscillates
uniformly (when \(k_t \neq 0\))

This structure functions as the geometric representation of force
``fields.'' A spherical wave emitted from a point source spreads over
the entire \(S^3\), extending the range of the force to the whole
subjective space.

\subsubsection{§4.5 CP Symmetry}\label{cp-symmetry}

In Table A of {[}4{]}, spherical waves with \(k_t = 0\) (photon-type,
gluon-type) are distinguished from those with \(k_t \neq 0\)
(\(W^{\pm}\)-type). By Proposition 4.2:

\textbf{Proposition 4.3}: CP symmetry can be violated only in
interactions mediated by spherical waves with \(k_t \neq 0\).
Interactions mediated by spherical waves with \(k_t = 0\) are
automatically CP-symmetric.

In Table A of {[}4{]}, the only gauge boson with \(k_t \neq 0\) is
\(W^{\pm}\), which is consistent with the fact that CP violation is
observed only in weak interactions. The naming conventions such as
``photon-type'' and ``\(W^{\pm}\)-type'' here indicate correspondence
with the interpretation layer of {[}4{]}; Proposition 4.3 itself relies
solely on the geometric condition of whether \(k_t = 0\).

\begin{center}\rule{0.5\linewidth}{0.5pt}\end{center}

\subsection{\texorpdfstring{§5 Standing Waves (\(\kappa = 1\),
\(\kappa_s = 1\),
\(\kappa_t = 0\))}{§5 Standing Waves (\textbackslash kappa = 1, \textbackslash kappa\_s = 1, \textbackslash kappa\_t = 0)}}\label{standing-waves-kappa-1-kappa_s-1-kappa_t-0}

\subsubsection{§5.1 Standing Oscillation Along One Spatial
Axis}\label{standing-oscillation-along-one-spatial-axis}

For \(\kappa = 1\) with \(\kappa_s = 1\), \(\kappa_t = 0\), the wave
vector is \(\mathbf{k} = (k_x, 0, 0, 0)\):

\[\psi(\mathbf{x} + \mathbf{e}_1) = e^{ik_x} \cdot \psi(\mathbf{x}) \quad (\text{phase change in the $x$ direction})\]

\[\psi(\mathbf{x} + \mathbf{e}_s) = \psi(\mathbf{x}) \quad (s \neq 1,\; \text{invariant in other directions})\]

\[\psi(\mathbf{x}, t + 1) = \psi(\mathbf{x}, t) \quad (\text{time-invariant})\]

This wave has an amplitude distribution along one spatial axis and is a
standing structure that does not change in time.

\subsubsection{§5.2 Scalar Character}\label{scalar-character}

The spatial axes \(k_1, k_2, k_3\) are equivalent by the
\(\mathrm{SO}(3)\) symmetry of \(S^3\). While
\(\mathbf{k} = (k_x, 0, 0, 0)\) selects the \(x\)-axis,
\(\mathbf{k} = (0, k_y, 0, 0)\) and \(\mathbf{k} = (0, 0, k_z, 0)\)
exist equally.

The average over the three equivalent configurations is isotropic, which
is the condition for a scalar field.

\textbf{Proposition 5.1}: Waves with \(\kappa = 1\), \(\kappa_t = 0\)
define a scalar quantity by virtue of the equivalence of three
directions under \(\mathrm{SO}(3)\) symmetry.

\subsubsection{§5.3 Filling of the Subjective Space by the Ground
Mode}\label{filling-of-the-subjective-space-by-the-ground-mode}

By the standing wave condition from {[}1{]}, the ground mode is a wave
with wavelength \(\lambda = 2\pi R\) that extends over the entire
\(S^3\). When a standing wave with \(\kappa = 1\) takes the ground mode:

\begin{itemize}
\tightlist
\item
  It fills the entire \(S^3\) with one wavelength
\item
  It is time-invariant (\(k_t = 0\))
\item
  It has phase structure along only one spatial axis
\item
  It is scalar by \(\mathrm{SO}(3)\) averaging
\end{itemize}

This is the description of the simplest non-trivial standing wave
filling a closed space and functions as the geometric representation of
vacuum structure.

\subsubsection{\texorpdfstring{§5.4 Uniqueness of
\(\kappa = 1\)}{§5.4 Uniqueness of \textbackslash kappa = 1}}\label{uniqueness-of-kappa-1}

\(\kappa = 0\) is trivial (constant function) within \(\mathbb{Z}^4\)
and has no spatial phase structure. \(\kappa = 2\) inherently possesses
surface orientation (rotational structure) as shown in §6 and is not
scalar.

\textbf{Proposition 5.2}: The minimal structure for a spatially
non-trivial and scalar (non-rotational) shape-invariant wave is
\(\kappa = 1\), which is the unique scalar mode.

\subsubsection{§5.5 Comparison with the Spherical Wave
Mode}\label{comparison-with-the-spherical-wave-mode}

The spherical wave with \(\kappa = 0\) (§4.4) and the standing wave with
\(\kappa = 1\) both fill the entire \(S^3\) as ground modes, but they
are qualitatively different:

{\def\LTcaptype{none} % do not increment counter
\begin{longtable}[]{@{}lll@{}}
\toprule\noalign{}
& Spherical wave (\(\kappa = 0\)) & Standing wave (\(\kappa = 1\)) \\
\midrule\noalign{}
\endhead
\bottomrule\noalign{}
\endlastfoot
Spatial phase structure & None & Exists along one axis \\
Spin & 0 or higher (depends on extended axes) & 0 (scalar) \\
Propagation & Can spread spherically & Standing (does not propagate) \\
Role & Force mediation (field) & Vacuum structure (ground state) \\
\end{longtable}
}

\begin{center}\rule{0.5\linewidth}{0.5pt}\end{center}

\subsection{§6 Geometric Origin of
Chirality}\label{geometric-origin-of-chirality}

\subsubsection{§6.1 Two-Axis Plane and Winding
Direction}\label{two-axis-plane-and-winding-direction}

Waves with \(\kappa = 2\) possess a plane defined by two nonzero axes.
On this plane, the winding direction of the phase field (clockwise /
counterclockwise) is defined. For the winding direction to be physically
distinguishable, the orientation of the plane must be fixed externally.

\subsubsection{\texorpdfstring{§6.2 (Space, \(t\)) Plane: Definition of
Chirality}{§6.2 (Space, t) Plane: Definition of Chirality}}\label{space-t-plane-definition-of-chirality}

For \(\mathbf{k} = (k_x, 0, 0, k_t)\), on the \((x, t)\) plane:

\begin{itemize}
\tightlist
\item
  The \(t\)-axis has its future direction fixed (arrow of time)
\item
  The \(x\)-axis is a spatial coordinate (selectable by
  \(\mathrm{SO}(3)\) rotation, but fixed once selected)
\end{itemize}

Since the direction of \(t\) is fixed, the two winding directions of the
phase on the \((x, t)\) plane are physically distinguishable:

\[\chi = \mathrm{sign}(k_x \cdot k_t) = \begin{cases} +1 & (\text{right-handed}) \\ -1 & (\text{left-handed}) \end{cases}\]

\textbf{Definition 6.1}: For a wave with \(\kappa = 2\),
\(\kappa_s = 1\), \(\kappa_t = 1\), the sign product
\(\chi = \mathrm{sign}(k_{s_0} \cdot k_t)\) of the nonzero spatial
component \(k_{s_0}\) (the unique spatial component with
\(k_{s_0} \neq 0\)) and the temporal component \(k_t\) is defined as the
\textbf{chirality}.

\textbf{Proposition 6.1}: Waves with \(\kappa = 2\) possessing a (space,
\(t\)) plane have two states, left-handed and right-handed, by virtue of
the fixed direction of the \(t\)-axis. This distinction cannot be
removed by continuous spacetime transformations.

\subsubsection{§6.3 (Space, Space) Plane: Degeneracy of
Chirality}\label{space-space-plane-degeneracy-of-chirality}

For \(\mathbf{k} = (k_x, k_y, 0, 0)\), on the \((x, y)\) plane:

\begin{itemize}
\tightlist
\item
  Both the \(x\) and \(y\) axes are spatial coordinates
\item
  An \(\mathrm{SO}(2)\) rotation \(R(\pi)\) can reverse the orientation
  of the plane
\end{itemize}

Since the orientation of the \((x, y)\) plane can be freely changed by
continuous rotation, clockwise and counterclockwise are physically
indistinguishable.

\textbf{Proposition 6.2}: Waves with \(\kappa = 2\) possessing only
(space, space) planes have degenerate chirality; left-handed and
right-handed are physically indistinguishable.

\subsubsection{§6.4 Majorana Condition}\label{majorana-condition}

By the Feynman--Stuckelberg interpretation, an antiparticle is described
as a particle traveling backward in time. Waves with \(\kappa_t = 0\)
are invariant under time reversal:

\[\psi(\mathbf{x}, -t) = \psi(\mathbf{x}, t)\]

The forward-in-time wave and the backward-in-time wave describe the same
state.

\textbf{Proposition 6.3}: Waves with \(\kappa = 2\), \(\kappa_t = 0\)
are time-reversal symmetric, and the particle state and antiparticle
state are described by the same wave function (Majorana condition).

In Table A of {[}4{]}, the only fermions with \(\kappa_t = 0\) are
neutrinos. Proposition 6.3 geometrically suggests the possibility that
neutrinos are of Majorana type. The naming ``neutrino'' here indicates
correspondence with the interpretation layer of {[}4{]}; Proposition 6.3
itself is derived solely from the geometric conditions \(\kappa = 2\),
\(\kappa_t = 0\).

\subsubsection{§6.5 Consistency with Discrete
Transformations}\label{consistency-with-discrete-transformations}

We define four discrete transformations as operations on the wave
vector:

{\def\LTcaptype{none} % do not increment counter
\begin{longtable}[]{@{}
  >{\raggedright\arraybackslash}p{(\linewidth - 4\tabcolsep) * \real{0.2807}}
  >{\raggedright\arraybackslash}p{(\linewidth - 4\tabcolsep) * \real{0.3860}}
  >{\raggedright\arraybackslash}p{(\linewidth - 4\tabcolsep) * \real{0.3333}}@{}}
\toprule\noalign{}
\begin{minipage}[b]{\linewidth}\raggedright
Transformation
\end{minipage} & \begin{minipage}[b]{\linewidth}\raggedright
Action on wave vector
\end{minipage} & \begin{minipage}[b]{\linewidth}\raggedright
Chirality \(\chi\)
\end{minipage} \\
\midrule\noalign{}
\endhead
\bottomrule\noalign{}
\endlastfoot
Parity \(P\) & \(k_x \to -k_x\), \(k_t\) unchanged & Reversed \\
Time reversal \(T\) & \(k_t \to -k_t\), \(k_x\) unchanged & Reversed \\
Charge conjugation \(C\) & \((k_x, k_t) \to (-k_x, -k_t)\) &
\textbf{Unchanged} \\
\(CPT\) & All components reversed then restored & \textbf{Unchanged} \\
\end{longtable}
}

\textbf{Proposition 6.4}: The \(CPT\) transformation is the reversal of
all components of the wave vector and holds automatically as a lattice
symmetry of \(\mathbb{Z}^4\).

\begin{center}\rule{0.5\linewidth}{0.5pt}\end{center}

\subsection{§7 Geometric Origin of the Pauli Exclusion
Principle}\label{geometric-origin-of-the-pauli-exclusion-principle}

\subsubsection{§7.1 Signed Area}\label{signed-area}

The two-axis plane defined by a wave with \(\kappa = 2\) possesses an
orientation. The orientation of a two-dimensional plane has \textbf{only
two possibilities: positive (front) and negative (back)}.

Consider the winding number \(n \in \mathbb{Z}\) of the phase field. By
Proposition 9.1 of {[}2{]} (Nyquist stability condition), the phase
change along each axis is restricted to \(|\phi_i| \leq \pi/2\). The
total phase change for one cycle on the two-axis plane is restricted to
\(|\phi_1| + |\phi_2| \leq \pi\), and a phase field with winding number
\(n\) requires a phase change of \(2\pi |n|\) per cycle. From
\(|2\pi n| \leq \pi\), we get \(|n| \leq 1/2\), and combined with the
condition \(n \in \mathbb{Z}\), \(n\) is limited to
\(n \in \{-1, 0, +1\}\) (\(n = 0\) is trivial). Therefore, fundamental
shape-invariant waves have \(n = \pm 1\), giving:

\[n = +1 \quad \Rightarrow \quad \text{signed area} = +A\]

\[n = -1 \quad \Rightarrow \quad \text{signed area} = -A\]

\subsubsection{§7.2 Maximum Coexistence Number of Identical
States}\label{maximum-coexistence-number-of-identical-states}

Consider a situation where multiple waves with \(\kappa = 2\) sharing
the same wave vector \(\mathbf{k}\) coexist in the same lattice region:

\begin{itemize}
\tightlist
\item
  First wave: signed area \(+A\)
\item
  Second wave: signed area \(-A\) (opposite orientation)
\item
  Third wave: Only two orientations, \(+A\) and \(-A\), exist for the
  two-dimensional plane. Both are already occupied, and no independent
  third state exists.
\end{itemize}

\textbf{Proposition 7.1} (Exclusion principle): Shape-invariant waves
with \(\kappa = 2\) sharing the same wave vector can coexist in the same
lattice region up to a maximum of two, because only two orientational
states (\(\pm\)) exist for the two-axis plane.

\subsubsection{\texorpdfstring{§7.3 Cases with
\(\kappa \neq 2\)}{§7.3 Cases with \textbackslash kappa \textbackslash neq 2}}\label{cases-with-kappa-neq-2}

\(\kappa = 0\): No nonzero axes exist in space \(\to\) no two-axis plane
can be formed \(\to\) no orientational constraint.

\(\kappa = 1\): Only one axis \(\to\) one axis short of forming a plane
\(\to\) no orientational constraint.

\textbf{Proposition 7.2}: Waves with \(\kappa \neq 2\) have no signed
area, and arbitrarily many waves with the same wave vector can coexist.

\subsubsection{§7.4 Geometric Correspondence of
Spin-Statistics}\label{geometric-correspondence-of-spin-statistics}

Combining Propositions 7.1 and 7.2:

{\def\LTcaptype{none} % do not increment counter
\begin{longtable}[]{@{}llll@{}}
\toprule\noalign{}
\(\kappa\) & Signed area & Maximum coexistence number & Statistics \\
\midrule\noalign{}
\endhead
\bottomrule\noalign{}
\endlastfoot
2 & Present (\(\pm\)) & 2 & Fermi-type \\
0, 1 & Absent & Unlimited & Bose-type \\
\end{longtable}
}

The geometric correspondence of the spin-statistics theorem --
half-integer spin obeys Fermi statistics, integer spin obeys Bose
statistics -- is obtained from the fact that ``a two-dimensional plane
has only two orientations.''

\subsubsection{§7.5 Pair Condensation}\label{pair-condensation}

A pair consisting of a wave with signed area \(+A\) and a wave with
\(-A\) has zero net signed area:

\[(+A) + (-A) = 0\]

A pair with zero area is not subject to the constraint of Proposition
7.1.

\textbf{Proposition 7.3}: A pair of waves with \(\kappa = 2\) (a
combination of areas \(\pm A\)) has zero net signed area, and as a pair
unit, is not subject to the coexistence constraint.

This corresponds to the geometric description of Cooper pair
condensation in superconductors (the phenomenon where pairs of fermions
exhibit bosonic behavior).

\begin{center}\rule{0.5\linewidth}{0.5pt}\end{center}

\subsection{§8 Child-Space Central Projection and
Acceleration}\label{child-space-central-projection-and-acceleration}

\subsubsection{§8.1 Uniform Motion and
Acceleration}\label{uniform-motion-and-acceleration}

We restate the discrete flow on the tangent bundle
\(\mathbb{Z}^4 \times \mathbb{Z}^4\) introduced in §8 of {[}2{]}:

\[\mathbf{k}_{t+1} = \mathbf{k}_t + \Delta\mathbf{k}_t \quad (\text{position update})\]

\[\Delta\mathbf{k}_{t+1} = \Delta\mathbf{k}_t + \Delta^2\mathbf{k}(\mathbf{k}_t) \quad (\text{velocity update})\]

The first equation describes uniform motion, and \(\Delta^2\mathbf{k}\)
in the second equation corresponds to acceleration. When
\(\Delta^2\mathbf{k} = 0\), the motion is uniform and rectilinear; when
\(\Delta^2\mathbf{k} \neq 0\), it is accelerated motion.

\subsubsection{§8.2 Global Acceleration: Curvature of the Subjective
Space}\label{global-acceleration-curvature-of-the-subjective-space}

As shown in {[}1{]}, shape-invariant standing waves exist on \(S^3\),
and the curvature radius \(R\) of \(S^3\) is not an externally given
constant but a quantity self-consistently determined from the
ground-mode structure of the standing waves present in the system
({[}1{]} §3.2: the standing wave condition \(2\pi R = n \cdot \lambda\)
determines \(R\) a posteriori). Geodesics (great circles) on \(S^3\)
follow the curvature. When these geodesics are centrally projected onto
\(\mathbb{Z}^4\), they appear as curves, and their curvature gives rise
to the global acceleration:

\[\Delta^2\mathbf{k}_{\text{global}} = f(R, \mathbf{k}_t)\]

Here \(f\) is a function depending on the curvature scalar \(R\) of
\(S^3\) and the current position \(\mathbf{k}_t\), determined from the
Jacobian of the central projection.

\subsubsection{§8.3 Local Acceleration: Child-Space Central
Projection}\label{local-acceleration-child-space-central-projection}

We utilize the central projection structure shown in §10 of {[}2{]}. In
the central projection \(\pi : S^5 \to \mathbb{Z}^4\) from the parent
space \(S^5(R)\) to the child space \(\mathbb{Z}^4\):

\begin{itemize}
\tightlist
\item
  Wave trajectories in the parent space are geodesics (segments of great
  circles)
\item
  The central projection is a fractional (nonlinear) map
\item
  In the child space, geodesics appear as curves
\end{itemize}

The local acceleration is determined by the nonlinearity of the central
projection:

\[\Delta^2\mathbf{k}_{\text{local}}(\mathbf{k}_t) = D^2\pi(\mathbf{k}_t) \cdot \Delta\mathbf{k}_t\]

Here \(D^2\pi\) is the Hessian matrix of the central projection \(\pi\),
evaluated discretely on \(\mathbb{Z}^4\).

\textbf{Proposition 8.1}: The second derivative \(D^2\pi\) of the
child-space central projection geometrically determines the acceleration
of shape-invariant wave trajectories. This acceleration is not a force
imposed externally but arises intrinsically from the curvature of the
projection.

\subsubsection{§8.4 Geometric Unification}\label{geometric-unification}

Global acceleration (curvature of \(S^3\)) and local acceleration
(nonlinearity of the central projection) are descriptions of the same
geometric mechanism at different scales:

\begin{itemize}
\tightlist
\item
  \textbf{Global}: The energy distribution of many waves determines the
  curvature radius \(R\) of \(S^3\) \(\to\) \(R\) affects all wave
  trajectories (integral effect)
\item
  \textbf{Local}: The Hessian matrix \(D^2\pi\) of the projection in the
  neighborhood of an individual wave produces local trajectory
  deflection (differential effect)
\end{itemize}

Both are connected as two aspects of a single mechanism because the
central projection \(\pi: S^5(R) \to \mathbb{Z}^4\) is a map that
parametrically depends on the curvature radius \(R\), and \(R\) itself
is self-consistently determined from the distribution of waves (§8.2).
The determination of \(R\) (global) and the evaluation of \(D^2\pi\) at
each point (local) are different aspects of the same \(R\)-parameterized
map \(\pi\).

\textbf{Proposition 8.2}: Acceleration is not an additional term in the
wave equation but is intrinsically derived from the geometry of the
central projection. Global curvature (determination of \(R\)) and local
projection nonlinearity (point-by-point evaluation of \(D^2\pi\)) are
manifestations at different scales of the single \(R\)-parameterized
central projection map.

\begin{center}\rule{0.5\linewidth}{0.5pt}\end{center}

\subsection{§9 Summary}\label{summary}

We summarize the main results of this paper:

\begin{enumerate}
\def\labelenumi{\arabic{enumi}.}
\item
  \textbf{Mode classification}: Shape-invariant waves are classified
  into three types by the number of nonzero components \(\kappa\) of the
  wave vector: \(\kappa \in \{0, 1, 2\}\). No stable modes exist for
  \(\kappa \geq 3\) (§2).
\item
  \textbf{Direction-constrained waves} (\(\kappa = 2\)): They propagate
  in only one spatial direction, and the signed area of the two-axis
  plane exists. The choice of spatial axis gives rise to the generation
  structure (§3).
\item
  \textbf{Spherical waves} (\(\kappa = 0\)): They have no spatial
  directionality and propagate isotropically. When \(k_t = 0\), they are
  time-reversal symmetric and carry no intrinsic causal ordering. As a
  ground mode, they fill the entire subjective space and function as
  force fields (§4).
\item
  \textbf{Standing waves} (\(\kappa = 1\)): They are standing
  oscillations along one spatial axis and constitute the unique scalar
  mode by \(\mathrm{SO}(3)\) symmetry. As a ground mode, they fill the
  entire subjective space and function as vacuum structure (§5).
\item
  \textbf{Chirality}: In the (space, \(t\)) plane, left-handed and
  right-handed states are distinguished by the direction of the
  \(t\)-axis. In the (space, space) plane, chirality is degenerate and
  the Majorana condition is satisfied. CP violation is limited to
  interactions with \(k_t \neq 0\) (§6).
\item
  \textbf{Exclusion principle}: Because the orientation of a
  two-dimensional plane has only two states, at most two waves with
  \(\kappa = 2\) can coexist, while waves with \(\kappa \neq 2\) have no
  such restriction. Pairs (\(\pm A\)) have zero area and are bosonic
  (§7).
\item
  \textbf{Acceleration}: The second derivative of the child-space
  central projection geometrically determines acceleration, unifying
  global curvature and local nonlinearity (§8).
\end{enumerate}

The correspondence with Table A of {[}4{]} is shown below:

{\def\LTcaptype{none} % do not increment counter
\begin{longtable}[]{@{}
  >{\raggedright\arraybackslash}p{(\linewidth - 8\tabcolsep) * \real{0.1091}}
  >{\raggedright\arraybackslash}p{(\linewidth - 8\tabcolsep) * \real{0.1818}}
  >{\raggedright\arraybackslash}p{(\linewidth - 8\tabcolsep) * \real{0.4000}}
  >{\raggedright\arraybackslash}p{(\linewidth - 8\tabcolsep) * \real{0.1091}}
  >{\raggedright\arraybackslash}p{(\linewidth - 8\tabcolsep) * \real{0.2000}}@{}}
\toprule\noalign{}
\begin{minipage}[b]{\linewidth}\raggedright
Mode
\end{minipage} & \begin{minipage}[b]{\linewidth}\raggedright
\(\kappa\)
\end{minipage} & \begin{minipage}[b]{\linewidth}\raggedright
Table A correspondence
\end{minipage} & \begin{minipage}[b]{\linewidth}\raggedright
Spin
\end{minipage} & \begin{minipage}[b]{\linewidth}\raggedright
Statistics
\end{minipage} \\
\midrule\noalign{}
\endhead
\bottomrule\noalign{}
\endlastfoot
Direction-constrained wave & 2 (\(\kappa_s = 1, \kappa_t = 1\)) &
Charged fermions & 1/2 & Fermi \\
Planar standing wave & 2 (\(\kappa_s = 2, \kappa_t = 0\)) & Neutrinos &
1/2 & Fermi (Majorana) \\
Spherical wave & 0 & Gauge bosons (\(\gamma, Z, W^{\pm}, g, G\)) & 0, 1,
2 & Bose \\
Standing wave & 1 & Higgs \(H^0\) & 0 & Bose \\
\end{longtable}
}

All claims in this paper are stated as geometric facts, and the
correspondence with Table A is presented as a ``connectable
interpretation.'' Identification with physical particles does not affect
the validity of the propositions themselves.

\begin{center}\rule{0.5\linewidth}{0.5pt}\end{center}

\subsection{References}\label{references}

{[}1{]} Noriaki Kihara, ``Shape-Invariant Standing Waves on Closed
Spheres: Stability by Dimension and the Geometric Necessity of 3+1
Spacetime,'' 2026.

{[}2{]} Noriaki Kihara, ``Shape-Invariant Traveling Waves on the
4-Dimensional Integer Lattice: Existence Conditions, Conservation
Structures, and Self-Consistent Loops,'' 2026.

{[}3{]} Noriaki Kihara, ``A Three-Layer Model of the Universe via
Central Projection,'' Zenodo, 2025. DOI: 10.5281/zenodo.15083729.

{[}4{]} Noriaki Kihara, ``Set Structure and Combinatorial Properties of
the 6-Dimensional Hyperrectangle,'' Zenodo, 2025. DOI:
10.5281/zenodo.19688521.
\end{document}
